% docs/paper2/related-work.tex
% BODY of the related-work section for paper 2 (review point #21).
% No \section line and no preamble: main.tex supplies the heading and the label.
% Integration instructions: docs/paper2/related-work-integration.md
% Bibliography keys: all resolve in docs/paper2/references.bib.

\paragraph{Counterexample-guided inductive synthesis: the same loop, a different oracle.}
CEGIS \citep{solarlezama2006sketching,solarlezama2008thesis} alternates a \emph{learner}, which proposes a candidate consistent with the counterexamples seen so far, and a \emph{verifier}, which either certifies the candidate against a specification or returns a new counterexample. The guarantee is inherited entirely from the verifier: one that is \emph{complete relative to a specification} either proves the candidate on the whole input domain or produces a witness. SyGuS standardized the format \citep{alur2013sygus}; the oracle-guided theory \citep{jha2010oracle,jha2017theory} makes what is synthesizable a function of which oracles are available, with programming by example \citep{gulwani2011flashfill,gulwani2017synthesis} as the bottom case, where generalization beyond the examples comes from a ranking function or a domain bias and not from a guarantee. Our synthesize--gate--refine pipeline \emph{is} a CEGIS loop --- LLM as learner, the gate's failing transitions as counterexamples, a refinement iteration as a CEGIS iteration --- with one load-bearing substitution: the counterexample oracle is random sampling of transitions, not a verifier. A verifier's ``no counterexample'' means \emph{there is none}; a sampling gate's means \emph{none was drawn in $N$ draws}, an event of probability exactly $(1-r)^N$ for a rule of rarity $r$ (Proposition~\ref{prop:gatemiss}), on which the loop is programming by example and the learner's prior decides everything off-sample. That weaker oracles buy weaker conclusions is not news to this literature and we do not re-derive it: it is the content of \citet{jha2017theory}, and the work that ran CEGIS against physical dynamics, where no complete verifier exists, substituted simulation and then recovered soundness by other means --- a bounded verifier, or an SMT check on the final candidate \citep{kapinski2014simulation,ravanbakhsh2019clf} --- or kept the verifier complete by pushing it into a theory solver \citep{abate2018cegist}. We add a closed-form law for how incomplete a sampling oracle is on a localized-rule failure, a proof that the miss event is unlearnable \emph{from the sample} by any learner (Proposition~\ref{prop:ident}), and a measurement of what a planner does with the accepted artifact. In the CEGIS sense our gate certifies one thing only --- the absence of counterexamples in a sample --- which is to say, nothing universal.

\paragraph{Statistical model checking.} SMC replaces exhaustive verification with simulation plus a statistical test and states what a sample size buys: acceptance sampling with Type-I/Type-II bounds \citep{younes2002acceptance}, Chernoff--Hoeffding complexity for an $(\epsilon,\delta)$ estimate of a property's probability \citep{herault2004apmc}, the black-box case \citep{sen2004blackbox}, and the standing difficulty that a rare property needs sample size scaling like the inverse of its probability \citep{legay2010smc}. Our gate is an SMC procedure for ``the model matches the plant to tolerance $\varepsilon$'' run without error bounds; we improve on none of that sample-complexity machinery, and \citet{herault2004apmc} is what would size a gate at a target confidence. SMC bounds the probability of a wrong verdict on a \emph{fixed} property, whereas we bound the probability that the sample never exhibits the phenomenon at all --- the verdict on the property actually tested is correct and uninformative --- and where SMC stops at the verdict we push it through a planner, whose loss is adversarially selected rather than proportional to the residual (Proposition~\ref{prop:playcost}).

\paragraph{Automata learning and active queries.} Passive data does not hand you the machine --- minimum-state automaton identification from given data is NP-hard \citep{gold1978complexity} --- whereas L$^*$ learns a regular language in polynomial time from membership queries \emph{plus an equivalence oracle} \citep{angluin1987learning}, which is what supplies a counterexample a random sample may never contain \citep{vaandrager2017model,settles2012active}. Proposition~\ref{prop:ident} is the continuous, hybrid-mode instance of that gap, and Angluin's decomposition reads our gate in one line: membership queries drawn at random, no equivalence oracle, which is why more refinement iterations cannot complete the loop. We inherit no query-complexity result --- their target is finite and exactly identifiable under exact queries, ours a program over a continuous space with tolerance $\varepsilon$ and a stochastic query distribution --- which is why we prove a probability law instead.

\paragraph{Falsification and counterexample-guided refinement.} Falsification searches for inputs driving a hybrid or cyber-physical system to violate a specification \citep{annpureddy2011staliro,corso2021survey,yamagata2021falsification}; CEGAR refines an abstraction at the spurious counterexample \citep{clarke2000cegar,clarke2003hybridcegar,alur2006predicate}; counterexample-guided data augmentation feeds falsified inputs back as training data \citep{dreossi2018cegda,dreossi2019verifai}. The message is uniform, and we do not improve on it: directed search finds the rare violating input random sampling misses. Our gate is the passive dual --- uniform random testing whose blind spot is precisely the mode a falsifier would hunt for, with $(1-r)^N$ quantifying the gap --- and the contrast with \citet{dreossi2018cegda} is our thesis in one sentence: there the counterexample that fixes the model comes from a falsifier, here from the gate's own random draw, and on the miss event there is none. CEGAR is also the template for our distrust-region mitigation, differing in the bookkeeping: it refines an abstraction over-approximate by construction, so refinement preserves the guarantee, whereas the fence refines a \emph{trust set} around an artifact with no such property --- hence a measured collapse of the exploitation and, at the farthest two-dimensional knob, $7$ of $20$ episodes pinned at blind-level return (Section~\ref{sec:patch2d-mitigation}).

\paragraph{Property-based testing with adaptive generators.} The gate is random testing of a program against an oracle and $(1-r)^N$ is the standard rare-input coverage gap \citep{claessen2000quickcheck,rubinstein2017montecarlo}, but uniform generation is not where that field stopped: generators can be steered by search toward inputs maximizing a utility \citep{loscher2017targeted,loscher2018automating}, derived from the property so constrained inputs arise by construction \citep{lampropoulos2017luck}, or guided by coverage feedback \citep{padhye2019zest}. Ours sits at the uniform, non-adaptive end of that spectrum, and each of the three has an instantiation here that we do not implement: a generator targeting distance travelled toward a mode boundary, one derived from the contract so mode-entering transitions occur by construction, and coverage-guided rollout mutation whose signal is which branches of the synthesized program executed --- available because the model is code. What we add is not a generator but that the planner is the adversary which finds the untested branch at deployment.

\paragraph{Experiment design and active system identification.} What is recoverable is a property of the input signal and not only of the estimator: optimal input design \citep{mehra1974optimal}, identification-for-control and the revival of experiment design \citep{gevers2005identification}, identifiability under a given excitation \citep{ljung1999sysid}, and the learning-theoretic result that actively chosen inputs achieve sample complexities passive excitation cannot \citep{wagenmaker2020active}. This literature implies our prescription and we re-derive none of it: our gate is a passive, non-adaptive experiment, and Proposition~\ref{prop:ident} is its identifiability statement specialized to a hybrid mode and a code hypothesis class. There the design reduces \emph{parameter} variance in a fixed parametric family and the failure is graded; ours is discrete --- a rule is present or absent in the program --- and its consequence is read through a planner rather than an estimation-error norm. The active boundary probing we list as follow-up work (Section~\ref{sec:limitations}) is an experiment-design problem and we name it as one.

\paragraph{Identification of hybrid and piecewise-affine systems.} Recovering a piecewise-affine or hybrid system is a joint classification-and-regression problem --- the partition of the state-input space as well as the affine law in each region --- and it is classical: clustering \citep{ferraritrecate2003clustering}, mixed-integer programming \citep{roll2004identification}, bounded-error set membership \citep{bemporad2005boundederror}, algebraic geometry \citep{vidal2003algebraic}, the tutorial and mixed logical dynamical formulations \citep{paoletti2007hybrid,bemporad1999control}, and trace-mining of hybrid automata \citep{medhat2015mining,soto2021synthesis}. We contribute no identification algorithm; our question is what a \emph{sampling verifier certifies} when the mode is missed, and what a planner then does --- the case these methods assume away, since none forms a region containing no samples. One is already active: \citet{soto2019membership} synthesizes linear hybrid automata from membership queries, the closest existing form of the fix we prescribe. Proposition~\ref{prop:lipschitz} and Corollary~\ref{cor:locbudget} are also a statement about this literature's hypothesis classes --- an exactly localized disagreement is unavailable to a bounded-Lipschitz class at any volume --- which is why we take code rather than piecewise-affine with a fixed mode count.

\paragraph{Contact-rich and discontinuous dynamics.} Contact is the physical archetype of our instruments and this literature's finding runs in the same direction: smooth end-to-end learning is the wrong parameterization and a complementarity-structured implicit loss does far better \citep{pfrommer2021contactnets}, stiff contact makes the loss landscape pathological for deep learners as a fundamental matter \citep{parmar2021stiff}, and even standard analytic planar contact models are limited on real data \citep{fazeli2020limitations,fazeli2017contact}. \citet{parmar2021stiff} is the empirical shadow of Proposition~\ref{prop:lipschitz} and \citet{pfrommer2021contactnets} takes structurally our remedy --- put the discontinuity in the hypothesis class --- neither of which we improve upon. The difference is what fails: there contact is present in the data and the failures are accuracy failures; here the mode is absent and the artifact is \emph{certified}. Our repair-from-data result is the code analogue --- given contact transitions in the sample the synthesizer recovered the exact switching rule on the one-dimensional clamps, whereas the two smooth fits we ran spread the error instead (Table~\ref{tab:smooth}) --- and, like this literature, we stop short of contact-rich manipulation.

\paragraph{Hypothesis classes that can represent modes.} Classes that \emph{can} express mode structure are mature: mixtures of experts with a learned gate \citep{jacobs1991moe}, recurrent switching linear dynamical systems whose discrete state depends on the continuous one \citep{linderman2017rslds}, neural ODEs with learned event functions \citep{chen2021eventfn}, and neural hybrid automata that infer the number of modes and their transitions \citep{poli2021nha}. This family scopes our representational claim rather than being scoped by it: Proposition~\ref{prop:lipschitz} quantifies over \emph{Lipschitz} pairs, and a hard gate or a learned event function reintroduces exactly the unbounded local structure it requires, so it is not an argument that neural models cannot represent hybrid modes --- our MLP is a probe at $h = 8$ scale, not a baseline (Section~\ref{sec:limitations}), and Section~\ref{sec:smooth} measures a favourable learned event-function baseline from this family directly. What survives is class-relative on the representation side and learner-independent on the identifiability side: no architecture here can infer a mode from a sample that never enters it (Proposition~\ref{prop:ident}), and float-level exactness --- zero off-mode error at $\varepsilon = 10^{-9}$ (Table~\ref{tab:smooth}) --- is a property of any class containing the mode: the measured event-function baseline attains it on the 1D clamp exactly as code does, while the smooth fits cannot (Section~\ref{sec:smooth}).

\paragraph{Conformal and PAC guarantees for learned dynamics.} Conformal prediction gives any predictor distribution-free finite-sample coverage under exchangeability \citep{vovk2022alrw} and has been carried to dynamics and control: prediction regions over trajectories, planned against \citep{lindemann2023conformal}; learned monitors that reject unreliable violation predictions \citep{bortolussi2019npm}; PAC-Bayes bounds for policies across environments \citep{majumdar2021pacbayes}. We derive no better certificate, and the comparison sharpens our result rather than softening it: their coverage holds with respect to the \emph{calibration} distribution, which here is the gate distribution, and the planner is not exchangeable with it --- the region our certificate covers carries $1.9\%$ of the exploited planner's queries, and a region containing every step-$t$ level set carries $7.8\%$ (Section~\ref{sec:limitations}). Calibrating conformally on gate rollouts would attach to the same wrong artifact a marginal guarantee that is honest over the calibration distribution and silent about the planner's queries; a planner could treat the conformal sets as uncertainty and react to their width, but sets calibrated where the mode never fires carry no widening at the mode, so the failure is not an artifact of choosing a metric certificate. Conversely ours is a sup-norm statement with an explicit $\varepsilon + 2L\rho$ constant (Propositions~\ref{prop:coverage} and~\ref{prop:partition}), which conformal prediction does not give, and it needs a finite Lipschitz constant, which conformal prediction does not --- where ours is weakest, the mode's local constant being unbounded.

\paragraph{Runtime assurance, safety filters, and robust MPC.} Every guarantee in this family is anchored to something known a priori: a verified conservative baseline plus a switching rule \citep{seto1998simplex,sha2001simplicity}, a shield synthesized from a safety automaton \citep{alshiekh2018shielding}, recursive feasibility across iterations of a learned controller \citep{rosolia2018lmpc,hewing2020learningmpc}, a minimally invasive projection keeping a known backup viable \citep{wabersich2021safetyfilter}, or a bounded model-error description \citep{rawlings2017mpc}. Our failure is invisible to every anchor \emph{derived from the learned model or its training distribution}: the error is unbounded but confined to a region the gate accepted as error-free, so a robustness margin tuned to the off-mode residual does not see it and a safety filter built on the same accepted model inherits the blind spot. An anchor specified independently of the model --- a verified conservative baseline, a shield synthesized from a safety automaton --- does not inherit it; it simply pays its conservatism whether or not a mode was omitted. The distrust-region fence has the shape of a safety filter with the anchor moved --- assembled from deployment-time refutations rather than an a-priori safe set, correcting a certified-but-wrong model rather than an unsafe policy --- which is why we report it as a mitigation with measured behaviour, its two-dimensional failure included (Section~\ref{sec:patch2d-mitigation}), and not as a certificate.

\paragraph{Decision-aware model learning and objective mismatch.} Prediction loss is the wrong objective for a model that will be planned with, and this literature has formalized the right one: losses weighting error by its effect on the Bellman update \citep{farahmand2017vaml,farahmand2018itervaml}, the value equivalence principle --- models are equivalent if they induce the same Bellman updates over a set of functions and policies \citep{grimm2020vep} --- an existence proof at scale for training purely for planning utility \citep{schrittwieser2020muzero}, and the finding that prediction accuracy and control performance diverge while planners exploit model error \citep{lambert2020objective,janner2019mbpo,hafner2020dreamer}. Our headline is a value-equivalence statement: the gate tests \emph{prediction} equivalence on the gate distribution, and we exhibit an artifact prediction-equivalent to within $10^{-9}$ everywhere the gate looked that is value-\emph{anti}-equivalent, capturing essentially none of the attainable return at the knobs swept. Their constructive programme is to \emph{train} for value equivalence, which needs a value function or a policy at training time; here the model is written before the planner runs, so our prescription is a change of \emph{test distribution}, not of loss --- and their results concern approximation quality within a parametric class, whereas our failure is an all-or-nothing missing rule that no reweighting of a loss recovers from a sample that never touched it. Our contribution to that conversation is the \emph{verified-but-wrong localized} regime, which a bounded-Lipschitz model can represent only at the volume price of Corollary~\ref{cor:locbudget} and which neither learner we fit represents at all (Table~\ref{tab:smooth}), plus a closed-form acceptance-failure law confirmed at gate scale.

\paragraph{Code world models.} Code World Models \citep{lehrach2025cwm} and code-as-policy \citep{liang2023codeaspolicies,gao2023pal} supply the paradigm, and the companion paper \citep{aguilar2026verified} is the discrete study this paper extends: we transfer its provable core, show that its empirical residual (translation, not inference) does not transfer on the one-dimensional clamps, and add the localization budget separating code from bounded-Lipschitz hypotheses.
